\documentclass{article}
\usepackage{microtype,graphicx,booktabs,hyperref,amsmath,amssymb}
\newcommand{\theHalgorithm}{\arabic{algorithm}}
\usepackage[preprint]{icml2026}
\icmltitlerunning{Native Robot Timing and Benchmark Scene Alignment}
\makeatletter\gdef\icmlcorrespondingauthor@text{Author details pending}\makeatother
\begin{document}
\twocolumn[
\icmltitle{Native Robot Timing on a Preserved Laboratory Scene:\\Benchmark Alignment and an Audited Component Pilot}
\begin{icmlauthorlist}\icmlauthor{Research development draft}{draft}\end{icmlauthorlist}
\icmlaffiliation{draft}{Scoped experiment note; full learning-system evaluation remains open}
\vskip .3in]
\printAffiliationsAndNotice{}
\begin{abstract}
Time-constrained manipulation requires physical effects within the requested
window, rather than timely dispatch alone. We reuse the benchmark's preserved
A-r7 laboratory and execute three uninterrupted native trials for one fixed
contact-timing contract. A public cue requests bilateral finger contact five
seconds later, within 50 milliseconds. The original drive schedule is early;
dispatch at the target is late; anticipation using a previously frozen contact
delay meets the diagnostic window. All three complete the same bottle task.
We preserve raw contacts, states, continuous native clock records and actual
RGB frames. This pilot establishes a robot timing interface and strong baseline,
not neural learning, chemical calibration, hardware timing or benchmark-wide performance.
\end{abstract}

\section{Problem and benchmark alignment}
The proposed Paper B system generates, verifies, repairs and absorbs robot
experience under physical time contracts. The benchmark provides task programs,
scene identity, observation boundaries and independent event scoring. Experience
generation and policy absorption must use those same definitions.

The preserved A-r7 scene contains 56 rigid objects and nine articulations.
The pilot uses its FR3 A arm, original bottle, scale and carousel. Existing
V79/V80 evidence establishes bottle manipulation and released-boundary
reconstruction. It does not establish liquid restoration, chemical event
calibration or continuous reaction time through a cold rebuild.

Original task definitions are retained: C5 requires stopping addition within
two seconds of a color change; C6 specifies 120$\pm$10 seconds of stirring;
C13 couples a transient A event with a randomized A-to-B wait in the M1 track.
Our new B-T0 contact task is a separately labeled interface diagnostic. Its
public event message and native clock cannot count as no-clock M1 performance.

\section{Effective-event formulation}
Let independently measured cue and effect times be conservative intervals
$I_c=[\ell_c,u_c]$ and $I_e=[\ell_e,u_e]$ on the same identified clock.
For requested gap $D$ and tolerance $\epsilon$, the error interval is
\[ I_{\rm err}=[\ell_e-u_c-D,\ u_e-\ell_c-D]. \]
PASS requires $I_{\rm err}\subseteq[-\epsilon,\epsilon]$.
Overlap without containment is indeterminate; missing or censored evidence
is unavailable. A dispatch timestamp is distinct from an effective contact.
The onset bracket and its subsequent confirmation timestamp are both retained.
Three consecutive sampled contacts do not prove continuous contact between samples.

Clock IDs must match. Reconstructed workers form separate epochs unless an
independent continuous process clock and evolution are established. Renaming
the reference event or restarting a process cannot retrospectively repair the
original deadline. Host inference time is not subtracted from native time
without calibration.

\section{Prospective native protocol}
The Isaac Sim 5.1 / PhysX scene integrates at 240 Hz. We copy and verify 108
benchmark source/input files; the separate pilot freezes 115 execution inputs.
Each condition creates a fresh worker with its own uninterrupted task epoch.
Only joint drive targets and measured gravity compensation are used; task-time
root-pose and joint-state writes are prohibited.

The cue occurs after the ten-second approach prefix. The original skill closes
four seconds after this cue. Dispatch-at-target inserts 240 hold ticks and
issues closure at five seconds. Anticipation inserts 198 hold ticks and issues
closure at 4.825 seconds, using the previously measured 0.175-second contact
delay upper bound. No current-trial contact labels tune either schedule.
The original approach, closure shape and remaining bottle skill are retained.

The scorer reconstructs the first run of three bilateral finger contacts from
raw PhysX impulses, rechecks continuous integration and compares every drive
target with the frozen command source. It separately recomputes manipulation
conditions from observed states and contacts. RGB capture does not advance
physics; 84 native frames per condition form three seven-second clips.
Post-run video annotations are presentation data, never policy inputs.

\section{Component result and limits}
\begin{table}[t]\centering\small
\caption{One fixed instance, three fresh trials. Times are seconds after the cue.}
\begin{tabular}{lrrl}\toprule
Condition & Dispatch & Effective contact & Verdict\\\midrule
Original & 4.000 & 4.171--4.175 & EARLY \\
Dispatch at target & 5.000 & 5.171--5.175 & LATE \\
Anticipation & 4.825 & 4.996--5.000 & PASS \\
\bottomrule\end{tabular}\end{table}
The three complete trajectories contain 40,719 measured task samples.
All finish the bottle manipulation and pass the existing mechanical conditions.
The timing difference follows deterministic actuator latency and scheduling;
it does not demonstrate adaptation or a learned improvement. There is one
independent task instance, with no statistical claim across seeds.

\begin{figure}[t]\centering
\includegraphics[width=\columnwidth]{figures/native_robot_timing_compact.pdf}
\caption{Squares: dispatch. Circles/bars: sampled physical contact. The shaded band is the requested window.}
\end{figure}

Formal chemical and training admission remain false. No Astra call, new
DP/RL/OPD weights, real hardware rollout or chemical outcome is credited.
Models do not run in this pilot, so it cannot establish physical progress
during planner waits or real inference deadlines. Controlled worker exit
preserves evidence but does not close graceful SDK shutdown acceptance.

\section{Next scientific tests}
The learning study must compare a shared history-conditioned policy against
strong timing rules under identical information and experience budgets.
Separate successful, failed and recovered data; compare offline expert
imitation with student-state relabeling; test context reset/shuffle and
changed dynamics. Split task instances before learning and separately hold
out state variation and goal/order/dependency variation.

Chemical extension requires independently calibrated flow-stop, stirring,
heater/temperature and endpoint events. Active processes must continue through
planning delays; long-task restoration must preserve the original reference
and account for every discontinuity. Until those gates pass, this component
supports interface development rather than the full method-paper claim.
\end{document}
